% \chapter{事例重建与选择策略}
\chapter{Event Reconstruction and Selection}

\section{Overview of the Analysis Strategy}
\label{sec:analysis-overview}

% [TODO] 简要说明分析策略的整体逻辑：从探测器信息出发，逐级重建物理对象，最终构建事例级候选并给出筛选后的事例样本。

\textbf{[To be filled: A short overview of the analysis strategy. The chapter proceeds from low-level reconstruction (primary vertex, tracks) through intermediate physics objects (muons, $J/\psi$, $\phi$) to event-level candidates, ending with the selected sample of 48,829 events that enters the signal extraction.]}

\section{Primary Vertex and Track Reconstruction}
\label{sec:pv-track-reco}

In CMS event reconstruction, primary vertex (PV) reconstruction is broadly divided into two steps: vertex finding and vertex fitting. The former groups charged tracks produced in the same bunch crossing according to their spatial positions, forming several vertex candidates. The latter estimates the vertex position and its covariance for each candidate track set, and provides the relevant track parameters and covariances after applying the vertex-position constraint, together with the weight of each track in the vertex fit. Through this vertex reconstruction procedure, the spatial positions of collision vertices, their positions relative to the beam line, and the association between each track and each vertex are estimated, providing essential spatial and association information for subsequent physics-object reconstruction and event selection.

The reconstruction procedure first selects, from the charged-track collection and using beam-spot information, the charged tracks suitable for primary vertex reconstruction. The default selection requirements are shown in Table~\ref{tab:pv-reco-track-selection}. These requirements remove tracks with poor reconstruction quality, large impact parameters relative to the beam line, large measurement uncertainties, or trajectories outside the main geometrical acceptance of the tracking detector, thereby improving the stability of the vertex-finding procedure.

\begin{table}
  \centering
  \begin{threeparttable}[c]
    \caption{Default charged-track selection requirements for primary vertex reconstruction}
    \label{tab:pv-reco-track-selection}
    \begin{tabular}{ll}
      \toprule
      Quantity & Default requirement \\
      \midrule
      Normalized track-fit $\chi^2$\tnote{a} & $< 10$ \\
      Transverse impact-parameter significance $|d_0|/\sigma_{d_0}$ & $< 4$ \\
      Transverse impact-parameter uncertainty $\sigma_{d_0}$ & $< \SI{1.0}{cm}$ \\
      Longitudinal impact-parameter uncertainty $\sigma_{d_z}$ & $< \SI{1.0}{cm}$ \\
      Transverse momentum $p_{\mathrm{T}}$ & $> \SI{0}{GeV}$ \\
      Pseudorapidity $\eta$ & $|\eta| < 2.4$ \\
      Number of pixel-detector hit layers & $\ge 2$ \\
      Number of silicon-tracker hit layers & $\ge 5$ \\
      Track-quality flag & Any quality flag accepted \\
      \bottomrule
    \end{tabular}
    \begin{tablenotes}
      \item [a] The normalized $\chi^2$ is obtained from the track Kalman-filter fit and reflects the overall quality of track reconstruction.
    \end{tablenotes}
  \end{threeparttable}
\end{table}

The preselected tracks are then clustered according to their positions along the beam direction. The $z$ coordinate used for clustering is the $z$ coordinate $z_{\mathrm{IP}}$ of the impact point (IP), defined as the point of closest approach of the track to the beam line.

Track clustering must also account for the uncertainty in $z_{\mathrm{IP}}$. This uncertainty receives contributions from the intrinsic $z$-position uncertainty $\sigma_{d_z}$ provided by the track reconstruction algorithm, as well as from the transverse width of the beam spot, which induces an additional $z$-direction uncertainty $\sigma_{z,\mathrm{BS}}$ in the track impact-point position. In the default configuration, the track-clustering algorithm also accounts for the intrinsic scale of the primary-vertex position by adding an intrinsic vertex-size term. This intrinsic vertex scale is:
\begin{equation}
  \sigma_{\mathrm{vtx}} = \SI{0.006}{cm}.
\end{equation}

This gives the effective $z$-direction position uncertainty $\sigma_{z,\mathrm{eff}}$:
\begin{equation}
  \sigma_{z,\mathrm{eff}}^2
  =
  \sigma_{d_z}^2
  + \sigma_{z,\mathrm{BS}}^2
  + \sigma_{\mathrm{vtx}}^2
\end{equation}

During clustering, the weight assigned to each track is taken to be inversely proportional to $\sigma_{z,\mathrm{eff}}^2$. This reduces the influence of tracks with large $z$-position uncertainties on the clustering result and improves the stability and accuracy of the clustering procedure.

Track clustering is performed using a one-dimensional deterministic annealing (DA) algorithm in the $z$ direction~\cite{roseDeterministicAnnealingApproach1990,roseDeterministicAnnealingClustering1998}. The algorithm starts from a high effective temperature, allowing tracks to be associated with several vertex prototypes simultaneously through soft weights. The temperature is then lowered gradually, so that vertex candidates at different $z$ positions progressively separate and converge. The application of this algorithm to vertex finding in the CMS experiment is described in Ref.~\cite{Chabanat:2005zz}. In the default configuration, the annealing factor is $0.6$; during clustering, only vertex candidates within approximately $4\sigma\sqrt{T}$ of a track are considered; the annealing stopping temperature is $T_{\mathrm{stop}}=0.5$; and the temperature at which the vertex-splitting stage ends is $T_{\min}=2.0$. Intermediate vertex prototypes separated by less than $\SI{0.01}{cm}$ are merged. In addition, tracks with large transverse impact-parameter significance are downweighted in clustering, with a default downweighting scale of $d_0/\sigma_{d_0}=3$. \textbf{[Reference to be added: the above default parameters for DA vertex finding are taken from the vertex-reconstruction module configuration in the CMSSW source code (\texttt{github.com/cms-sw/cmssw}).]} These settings allow the algorithm to separate tracks from different pp interactions into multiple vertex candidates according to $z_{\mathrm{PCA}}$ in pileup environments, while reducing the impact of nonprimary or outlier tracks on the clustering result.

For each track cluster, a three-dimensional vertex fit is subsequently performed using the Adaptive Vertex Fitter (AVF). The AVF is an iterative weighted fitting method: at each iteration, the algorithm assigns weights to tracks according to their compatibility with the current vertex estimate. Tracks incompatible with the vertex position are gradually downweighted, while tracks compatible with the vertex retain high weights. The default geometric annealing parameters correspond to a $\chi^2$ cutoff scale of $2.5$. Thus, the AVF does not require a hard decision before the fit as to whether a given track belongs to a vertex; instead, it suppresses outlier tracks through soft weights during the fitting procedure. After the fit is completed, each vertex candidate must satisfy validity requirements and have a distance from the beam line in the transverse plane smaller than $\SI{1.0}{cm}$.

When multiple primary vertex candidates are present in an event, the vertices are ordered in descending order according to the sum of squared transverse momenta of their associated tracks. In the current implementation, the ranking variable is not the bare $\sum p_{\mathrm{T}}^2$, but rather
\begin{equation}
  S_{\mathrm{PV}}
  =
  \sum_{i \in \mathrm{PV}}
  \left[
    \max \left(
      p_{\mathrm{T},i} - \sigma_{p_{\mathrm{T}},i},\,0
    \right)
  \right]^2 ,
  \label{eq:pv-ranking}
\end{equation}
where $\sigma_{p_{\mathrm{T}},i}$ is the transverse-momentum uncertainty of the $i$th track. After sorting, the vertex with the largest $S_{\mathrm{PV}}$ is placed at index zero of the vertex collection and is usually regarded as the leading primary-vertex candidate for the hard-scattering process.

\begin{table}
  \centering
  \begin{threeparttable}[c]
    \caption{Main criteria for associating charged PF candidates with primary vertices}
    \label{tab:pv-association}
    \begin{tabular}{p{0.22\linewidth}p{0.68\linewidth}}
      \toprule
      Category & Criterion and physical interpretation \\
      \midrule
      Vertex-fit membership &
      If the track has already been used in the fit of a primary vertex, it is preferentially associated with that vertex. This is the strongest association and indicates that the track has already contributed to estimating the vertex position during vertex reconstruction.
      \\

      Prompt-like $d_z$\tnote{a} compatibility
      &
      If the track has not been explicitly used in a vertex fit, its longitudinal impact parameter with respect to each vertex is compared. By default, the closest compatible vertex must satisfy
      $|d_z| < \SI{0.1}{cm}$,
      $|d_z|/\sigma_{d_z,\mathrm{eff}} < 5$,
      and the track itself must have $\sigma_{d_z}<\SI{0.05}{cm}$. When these conditions are satisfied, the track is considered compatible with the vertex along the beam direction and consistent with the topology of a prompt charged particle.
      \\
      Heavy-flavor hadron decay track
      &
      If the track does not satisfy the prompt-like criteria, the algorithm further considers the possibility that it originates from the secondary decay of a long-lived particle such as a
      $B_s$ hadron. A nearby jet is then searched for, with the default requirements
      jet $p_{\mathrm{T}}>\SI{25}{GeV}$ and
      $\Delta R<0.5$ between the track and the jet. The track is subsequently required to satisfy
      $|d_z|<\SI{0.1}{cm}$ and $|d_{xy}|<\SI{0.1}{cm}$
      with respect to the candidate vertex, and its distance to the jet axis must be less than $\SI{0.07}{cm}$. If a vertex satisfying these criteria exists, the track is associated with that vertex, indicating that although the track is not prompt, it is compatible with a displaced decay topology inside a jet originating from that vertex.
      \\
      Unreconstructed primary-vertex case
      &
      If the track cannot be strongly associated with any reconstructed vertex, but has a small transverse impact parameter relative to the first vertex, namely $|d_{xy}|<\SI{0.01}{cm}$ and $|d_{xy}|/\sigma_{d_{xy}}<2$, the track is considered to possibly originate from a primary vertex that was not successfully reconstructed, and is assigned to the reconstructed vertex closest in the $z$ direction.
      \\
      Weak association to the nearest $z$ vertex
      &
      If none of the above criteria are satisfied but the track $d_z$ can still be computed with respect to each reconstructed vertex, the track is weakly associated with the vertex for which the $d_z$ significance is smallest. This step is used only as a fallback, indicating that among the available reconstructed vertices, this vertex is closest to the track in the $z$ direction.
      \\
      \bottomrule
    \end{tabular}
    \begin{tablenotes}
      \item [a] $d_z$ denotes the longitudinal impact parameter of the track with respect to the vertex, and $d_{xy}$ denotes the transverse impact parameter.
    \end{tablenotes}
  \end{threeparttable}
\end{table}

After primary vertex reconstruction is completed, the candidates reconstructed by the CMS particle-flow (PF) algorithm are associated with reconstructed primary vertices. For each charged PF candidate, the algorithm first checks whether its charged track has already been used in a vertex fit and has a nonzero fit weight in that vertex. If this condition is satisfied, the vertex is considered the most reliable origin vertex for the particle. This is because the track has already contributed to the vertex-position estimate during the vertex-fitting stage and has high compatibility with that vertex. If a track cannot be explicitly assigned through its vertex-fit weight, the hierarchical criteria listed in Table~\ref{tab:pv-association} are applied. This hierarchy reflects physical information from strongest to weakest: vertex-fit membership is used first, followed by prompt-like $d_z$ compatibility, then displaced tracks from heavy-flavor hadron decays, and finally the nearest $z$ distance as a weak association.

Charged-particle tracks in the CMS silicon tracker are reconstructed using an iterative Combinatorial Track Finder (CTF) based on the Kalman filter method~\cite{Fruhwirth:1987:Kalman,CMS_TDR_CH6_TRACKING}. The reconstruction sequence proceeds through hit clustering, seed generation, pattern recognition (track finding), ambiguity resolution (removal of redundant tracks sharing too many hits), and final track fitting.

Raw signals are first clustered in the pixel and strip sensors. In the pixel detector, charge sharing among adjacent pixels provides two-dimensional position information with a spatial resolution of about 10~$\mu$m in the bending plane. In the strip detector, signals on adjacent strips above threshold are merged into one-dimensional hit clusters. After clustering, track reconstruction begins with seed generation: combinations of two or three pixel hits, together with a hit in the first strip layer, are used to estimate initial five-dimensional helix parameters and their covariance matrix, providing starting values for the subsequent Kalman filter.

The CTF propagates each seed outward through successive tracker layers, using the Kalman filter to find hits compatible with the trajectory candidate and update the track parameters at each layer. When a candidate can be propagated through multiple hit assignments, the algorithm maintains several trajectory candidates in parallel and prunes them at each layer based on $\chi^2$ compatibility with the assigned hits. After traversing all tracker layers, a second Kalman filter (smoother) is applied from the outside in to obtain optimal parameter estimates at every measurement surface.

To improve the reconstruction efficiency across different track topologies, the CTF is run in multiple iterations. After each iteration, hits unambiguously assigned to high-quality tracks are removed from subsequent searches, reducing combinatorial complexity. In the Run~1 data processing, CMS employed six iterations, each targeting specific track topologies such as high-$p_{\mathrm{T}}$ prompt tracks, low-$p_{\mathrm{T}}$ tracks, and displaced tracks from long-lived particle decays.

If several reconstructed tracks share a majority of their hits, the ambiguity resolution step ranks tracks by number of hits and $\chi^2$, retains the best track, and removes redundant candidates. The final reconstructed tracks are classified into quality categories (loose, tight, high-purity) according to number of hits, $\chi^2$, and compatibility with the primary vertex; the high-purity category is used for physics analysis. For charged particles with $p_{\mathrm{T}} > 1$~GeV, the track reconstruction efficiency is approximately 95\%, the fake rate is below a few percent, and the transverse momentum resolution in the central region is about 1--2\% at $p_{\mathrm{T}} \sim 10$~GeV~\cite{CMS_TDR_CH6_TRACKING}.

\section{Muon Reconstruction and Selection}
\label{sec:muon-reco-selection}

Muon candidates in CMS are reconstructed by combining information from the silicon tracker and the muon spectrometer~\cite{CMS_MUON_PERF_13TEV,CMS_TDR_CH9_MUONS}. In the standard muon reconstruction sequence, tracks are first reconstructed independently in the inner silicon tracker and in the outer muon system. The latter procedure is referred to as standalone muon reconstruction: starting from track segments locally reconstructed in the drift tubes (DTs) and cathode strip chambers (CSCs), combined with hit information from the resistive plate chambers (RPCs), a Kalman filter builds tracks by propagating state vectors from the innermost muon station outward through successive stations. Energy loss, multiple scattering, and the non-uniform magnetic field in the iron return yoke are taken into account during propagation. A backward Kalman filter is then applied from the outside in, and the track parameters are defined at the innermost muon station, followed by an extrapolation to the nominal interaction point with a beam-spot constraint. The overall muon reconstruction efficiency is 95--99\% over $|\eta| < 2.4$, with slight dips at the geometrical boundaries between detector wheels and in the barrel--endcap transition region ($|\eta| \approx 0.25$, 0.8, and 1.2). Global muon reconstruction then matches standalone muon tracks to tracker tracks in the silicon tracker and performs a combined fit using hits from both the inner tracker and the muon system. Because it uses information from both detector systems, global muon reconstruction provides an improvement by nearly an order of magnitude in transverse momentum resolution compared to standalone reconstruction across all $p_{\mathrm{T}}$ ranges~\cite{CMS_TDR_CH9_MUONS}.

In addition to the outside-in reconstruction described above, CMS also uses an inside-out tracker muon reconstruction algorithm. This algorithm starts from reconstructed tracker tracks in the silicon tracker and extrapolates tracks satisfying
$p_{\mathrm{T}} > 0.5~\mathrm{GeV}$ and $p > 2.5~\mathrm{GeV}$
to the muon system. If the extrapolated track is matched to at least one DT or CSC segment, the track is labeled as a tracker muon. The matching is performed in the local coordinate system of the muon chamber and primarily uses the local $x$ coordinate, which has better measurement precision in the bending plane. A tracker track is considered matched to a muon segment when the position difference satisfies
$|\Delta x| < 3~\mathrm{cm}$
or the corresponding pull is less than 4. Tracker muons are particularly important for low-transverse-momentum muons and for detector regions with weaker geometrical coverage, where such muons may not form complete standalone or global muons.

In the High-Level Trigger system, muon reconstruction proceeds in two steps: Level-2 performs standalone track reconstruction in the muon spectrometer, seeded by Level-1 muon trigger candidates; Level-3 incorporates silicon-tracker information and performs a global fit reconstruction, seeded by Level-2 candidates~\cite{CMS_TDR_CH9_MUONS}. Under Run~2 data-taking conditions, the single-muon HLT path without isolation has a $p_{\mathrm{T}}$ threshold of approximately 19~GeV, while the typical dimuon threshold is 7~GeV.

This analysis uses the CMS soft muon working point as the muon identification criterion~\cite{CMS_MUON_PERF_13TEV}. This working point is optimized for low-transverse-momentum muons and is particularly suitable for $B$-physics and quarkonium-related final states. At the detector level, the soft muon working point is defined by the following requirements:

\begin{enumerate}
  \item The candidate must be reconstructed as a tracker muon, and its inner tracker track must satisfy the high-purity quality requirement. The tracker track must have hits in at least six silicon-tracker layers, including at least one pixel-detector hit;
  \item The segment matching for the tracker muon must be stringent: the pulls between the extrapolated tracker track and the matched muon segment in both the local $x$ and local $y$ coordinates must be less than 3, suppressing hadron punch-through and random matching backgrounds while retaining efficiency for low-transverse-momentum muons;
  \item The candidate must be loosely compatible with the primary vertex, with transverse and longitudinal impact parameters satisfying $|d_{xy}| < 0.3~\mathrm{cm}$ and $|d_z| < 20~\mathrm{cm}$, respectively.
\end{enumerate}

The central idea of the soft muon identification strategy is to rely primarily on the precise momentum and impact-parameter measurements provided by the inner silicon tracker, while using segment matching in the muon system to confirm the muon hypothesis. This achieves high muon identification efficiency in the low-transverse-momentum region.

On top of the soft muon identification, this analysis imposes the following kinematic selection requirements on muons:

\begin{itemize}
  \item \textbf{Pseudorapidity}: $|\eta| < 2.4$, ensuring that muons lie within the effective geometrical acceptance of the CMS silicon tracker and muon spectrometer.
  \item \textbf{Transverse momentum}: Region-dependent thresholds are used, based on the $\eta$ dependence of the trigger acceptance and reconstruction resolution of the muon detector:
    \begin{itemize}
      \item Barrel region ($|\eta| < 1.2$): $p_{\mathrm{T}} > 3.5~\mathrm{GeV}$;
      \item Endcap region ($1.2 \leq |\eta| < 2.4$): $p_{\mathrm{T}} > 2.5~\mathrm{GeV}$.
    \end{itemize}
\end{itemize}
This region-dependent $p_{\mathrm{T}}$ threshold design balances the requirements of trigger efficiency and analysis purity.

\section{\texorpdfstring{$J/\psi$}{J/psi} Candidate Reconstruction}
\label{sec:jpsi-reco}

The $J/\psi$ meson is reconstructed through its dimuon decay channel $J/\psi \to \mu^+\mu^-$. The branching fraction of this decay channel is approximately $5.96\%$~\cite{PDG2024}. Although this is smaller than the total branching fraction to light-hadron decay channels, about $87.7\%$, the two muons in the final state provide a clean experimental signature. With the excellent muon reconstruction and identification capabilities of CMS, this channel provides a high signal-to-background ratio and facilitates signal extraction in complex background environments.

$J/\psi$ candidates are reconstructed by selecting pairs of oppositely charged muons from the muons satisfying the selection requirements described in Section~\ref{sec:muon-reco-selection}. Each candidate is required to satisfy the following criteria:

\begin{enumerate}
  \item \textbf{Invariant-mass window}: The dimuon invariant mass $m_{\mu\mu}$ must satisfy $2.9 < m_{\mu\mu} < 3.3~\mathrm{GeV}$;
  \item \textbf{Transverse momentum and rapidity}: The transverse momentum of the $J/\psi$ candidate must satisfy $p_{\mathrm{T}}^{J/\psi} > 6.0~\mathrm{GeV}$, and its rapidity must satisfy $|y^{J/\psi}| < 2.5$;
  \item \textbf{Vertex-fit quality}: The two muon tracks must be fitted to a common decay vertex, and the vertex-fit probability must be greater than $1\%$.
\end{enumerate}

Among these selection requirements, the invariant-mass window balances the mass-peak resolution of the $J/\psi$ and the suppression of continuum background. The transverse-momentum and rapidity selections are based on the trigger efficiency and reconstruction performance of CMS. The vertex-fit quality requirement helps suppress combinatorial background and ensures that the two muons are spatially consistent with originating from the same decay vertex.

\section{\texorpdfstring{$\phi$}{phi} Candidate Reconstruction}
\label{sec:phi-reco}

The $\phi$ meson is reconstructed through its dominant decay channel $\phi \to K^+K^-$, which has a branching fraction of approximately $49.1\%$~\cite{PDG2024}. In the CMS experiment, the ability to distinguish $K^\pm$ from $\pi^\pm$ is limited for charged hadrons with transverse momenta at the $\mathrm{GeV}$ scale. Therefore, this analysis follows the strategy commonly used in CMS multiple-quarkonium analyses: all charged tracks that pass the basic quality selection and are not identified as muons are assigned the $K^\pm$ mass hypothesis~\cite{CMS_TRI_JPSI}.

$\phi$ candidates are reconstructed by selecting pairs of oppositely charged tracks from charged tracks satisfying the following requirements:

\begin{enumerate}
  \item \textbf{Track quality}: The charged tracks must satisfy the CMS high-purity quality requirement to ensure reliable track reconstruction.
  \item \textbf{Kinematic acceptance}: The track pseudorapidity must satisfy $|\eta| < 2.5$, and the transverse momentum must satisfy $p_{\mathrm{T}} > 2.0~\mathrm{GeV}$.
  \item \textbf{Invariant-mass window}: Under the $K^\pm$ mass hypothesis, the two-track invariant mass $m_{KK}$ must satisfy $0.99 < m_{KK} < 1.07~\mathrm{GeV}$, corresponding to an approximately $\pm 40~\mathrm{MeV}$ window centered on the PDG value of the $\phi$ mass, $1019.5~\mathrm{MeV}$~\cite{PDG2024}.
  \item \textbf{Transverse momentum}: The transverse momentum of the $\phi$ candidate must satisfy $p_{\mathrm{T}}^{\phi} > 4.0~\mathrm{GeV}$. This threshold is chosen based on the following physical consideration: the $\phi(s\bar{s})$ is a light-flavor meson and is abundantly produced in proton--proton collisions through soft QCD processes, mostly in the $p_{\mathrm{T}} \lesssim 1~\mathrm{GeV}$ region. Imposing a higher $p_{\mathrm{T}}$ threshold effectively suppresses combinatorial background from soft QCD processes and helps ensure that the reconstructed $\phi$ candidates originate from hard-scattering processes.
  \item \textbf{Vertex-fit quality}: The two charged tracks must be fitted to a common vertex, and the vertex-fit probability must be greater than $1\%$, suppressing combinatorial background from random track pairs.
\end{enumerate}

\section{\texorpdfstring{$J/\psi J/\psi\phi$}{J/psi J/psi phi} Candidate Reconstruction}
\label{sec:triplet-reco}

After reconstructing and selecting the individual final-state particles, namely muons, $J/\psi$ candidates, and $\phi$ candidates, this analysis constructs event-level $J/\psi + J/\psi + \phi$ candidate combinations:

\begin{enumerate}
  \item Two $J/\psi$ candidates satisfying the requirements described in Section~\ref{sec:jpsi-reco} are selected from the event, with the requirement that they do not share any muon tracks;
  \item One $\phi$ candidate satisfying the requirements described in Section~\ref{sec:phi-reco} is selected from the event, with the requirement that its two charged tracks do not overlap with the muon tracks used by the two $J/\psi$ candidates;
  \item The two $J/\psi$ candidates are fitted to a common vertex, with the vertex-fit probability required to be greater than $1\%$, ensuring that the two $J/\psi$ candidates are spatially consistent with originating from the same collision vertex~\cite{CMS_TRI_JPSI};
  \item The two $J/\psi$ candidates and the $\phi$ candidate are jointly fitted to the same vertex, and the fit is required to succeed, yielding a valid vertex, to further strengthen the signal hypothesis of simultaneous production of the three quarkonia.
\end{enumerate}

If multiple candidate combinations in an event satisfy the above requirements, the combination with the smallest $\chi^2$ from the common-vertex fit of the three final-state mesons is selected as the final candidate.

The hierarchical selection strategy above reflects the logical progression from individual physics objects to the complete event. First, the quality of single muons is ensured through the soft muon working point and kinematic selection. The $J/\psi$ candidates are then identified using the dimuon invariant mass and vertex fit, after which $\phi$ candidates are constructed in an analogous manner. Finally, the simultaneous production of the three quarkonia is tested at the event level through the common-vertex constraint on multiple mesons. This design suppresses background contributions from random combinations, continuum processes, and soft QCD processes step by step while retaining signal efficiency as much as possible.

\section{Trigger and Event-Level Selection}
\label{sec:trigger-event-selection}

The data used in this analysis are selected online by the CMS two-level trigger system. Candidate events in the analysis are required to pass the \texttt{HLT\_Dimuon0\_Jpsi3p5\_Muon2\_v} trigger path. The main event-selection requirements of this trigger path are:

\begin{itemize}
  \item At least three muons must be present in the event, each satisfying $p_{\mathrm{T}} > 2.0~\mathrm{GeV}$ and $|\eta| < 2.5$;
  \item At least one oppositely charged muon pair must satisfy $p_{\mathrm{T}} > 3.5~\mathrm{GeV}$;
  \item The invariant mass of the muon pair must satisfy $2.95 < m_{\mu\mu} < 3.25~\mathrm{GeV}$;
  \item The probability for the two muon tracks to be fitted to a common vertex must be greater than $0.5\%$.
\end{itemize}

The dimuon mass window of this trigger path is matched to the $J/\psi$ mass peak and is designed to efficiently select events containing $J/\psi \to \mu^+\mu^-$ decays.

\section{Selected Candidate Sample}
\label{sec:selected-candidate-sample}

Applying the full set of selection requirements described in Sections~\ref{sec:pv-track-reco}--\ref{sec:trigger-event-selection} to the Run~3 data sample yields a total of \textbf{48,829} $J/\psi J/\psi \phi$ candidate events. These events constitute the input to the signal extraction described in Chapter~5. The invariant mass distributions of the selected candidates, prior to any fit, are shown in Fig.~\ref{fig:prefit-mass-distributions}. The $J/\psi$ and $\phi$ resonance peaks are clearly visible above the combinatorial background.

\begin{figure}[htbp]
  \centering
  % \includegraphics[width=0.48\textwidth]{figures/chap04/prefit_mass_jpsi1.png}
  % \includegraphics[width=0.48\textwidth]{figures/chap04/prefit_mass_jpsi2.png}
  % \includegraphics[width=0.48\textwidth]{figures/chap04/prefit_mass_phi.png}
  \caption{Invariant mass distributions of the selected $J/\psi J/\psi \phi$ candidates before signal extraction: $m(J/\psi_1)$ (left), $m(J/\psi_2)$ (middle), and $m(\phi)$ (right).}
  \label{fig:prefit-mass-distributions}
\end{figure}
